\subsection{Russell's Paradox and Axioms}

\begin{theorem} \label{thm:russells-paradox} \hlt{Russell's Paradox}\\
For a condition $P(x)$ on sets, unrestricted comprehension asserts $(\exists R)(\forall x)\,[x\in R\Longleftrightarrow P(x)]$.
\begin{enumerate}[label=(\roman*)]
\setlength{\itemsep}{0pt}
\item No set $R$ satisfies $(\forall x)\,[x\in R\Longleftrightarrow x\notin x]$.
\item The unrestricted comprehension schema is inconsistent.
\end{enumerate}
\begin{proof}
\begin{enumerate}[label=(\roman*)]
\setlength{\itemsep}{0pt}
\item Suppose that such a set $R$ exists. Setting $x=R$ gives $R\in R\Longleftrightarrow R\notin R$. If $R\in R$, then $R\notin R$; if $R\notin R$, then $R\in R$.\\
Both cases are contradictory, so no such set exists.
\item Take $P(x)\defeq x\notin x$.
The corresponding instance asserts
$(\exists R)(\forall x)\,[x\in R\Longleftrightarrow x\notin x]$.
This contradicts (i), proving (ii).
\end{enumerate}
\hlt{Structural Commentary}: Specialize at $x=R$, then apply the matching comprehension instance.\\
\hlt{Key Insight}: Self-membership makes the defining condition assert its own negation.
\end{proof}
\end{theorem}

\begin{perspective} \label{per:zfc-axiom-role-map} \hlt{From Comprehension to Axioms}\\
The obstruction in Theorem \ref{thm:russells-paradox} explains the need for controlled set formation.\\
Separation forms $\{x\in B:P(x)\}$ only after a bounding set $B$ is available.
The standard ZFC axioms have the following informal roles.
\begin{enumerate}[label=(\roman*)]
\setlength{\itemsep}{0pt}
\item \hlt{Extensionality} makes membership determine set equality.
\item \hlt{Separation} extracts a condition-defined subset from an existing bounding set.
\item \hlt{Pairing} supplies $\{a,b\}$ from given sets $a$ and $b$, including $\{a\}$ when $a=b$.
\item \hlt{Union} collects into one set the elements of the members of a given set.
\item \hlt{Power set} supplies $\mathcal{P}(A)$, the set of all subsets of a given set $A$.
\item \hlt{Infinity} supplies an inductive set from which the natural numbers are constructed.
\item \hlt{Replacement} forms the image of a set under a definable functional relation.
\item \hlt{Foundation} requires nonempty sets to contain an $\in$-minimal member.
\item \hlt{Choice} supplies a choice function on a set-indexed family of nonempty sets.
\end{enumerate}
This perspective is a functional map, not a formal axiom list or a source of proof authority.\\
Formal statements and dependencies belong to the advanced set-theory section.
\end{perspective}
